% TOP_PROBLEM: {"id": 30004844, "title": "Local Smoothing for the Euclidean Wave Equation", "category_id": 8, "category": {"id": 8, "name": "analysis", "display_name": "Analysis", "description": "Limits, continuity, calculus, and function theory.", "slug": "analysis", "order_index": 8, "created_at": "2026-07-31T15:26:25.671Z"}, "status": "partially_solved"}
% FIRST_POSTED: 2026-10-26
% ATTEMPT_STATUS: unresolved
% SUBMISSION: {"kind": "research_attempt", "category": "top_problems", "problem_id": "30004844", "model": {"provider": "OpenAI", "version": "GPT 6 Astra Pro (exact release/build unknown)", "reasoning": "Unknown", "generated_on": "2026-10-26"}, "other_models": "No second model reported for exposition. Codex used for repository formatting; exact underlying release/build not recorded.", "tools": "Source inspection and TeX preparation reported in supplied manuscript; repository validation, LuaLaTeX and browser rendering during integration. No Lean execution.", "target": "Canonical local smoothing assertion for Euclidean waves in all spatial dimensions n>=2.", "scope": "Literature report only. Detailed exposition of externally produced OpenAI results and formalization scope; no new discovery or independent proof certification.", "human_contributions": "Submitter prompted preparation and expansion; no human mathematical authorship or independent proof verification reported.", "verification": "Source checks attributed to supplied report. Integration checks formatting and target identity, not all analytic/algorithmic proofs or Lean dependencies.", "reproducibility": "Supplied manuscript dated2026-10-26, pinned openai/math commit 7ef5c0cc5ca4b0e072b7a050fca09ab24bc6fa3a. Original prompts and browsing logs unavailable; initial/final research-prompt use not documented. Integration expands presentation/link macros and adds metadata without changing mathematical exposition."}
\documentclass[11pt,a4paper]{article}
\usepackage[T1]{fontenc}
\usepackage[utf8]{inputenc}
\usepackage{lmodern}
\usepackage[margin=25mm,headheight=14pt]{geometry}
\usepackage{amsmath,amssymb,amsthm,mathtools}
\usepackage{microtype,booktabs,array,longtable,enumitem}
\usepackage{xurl}
\usepackage[hidelinks]{hyperref}
\usepackage{fancyhdr}
\pagestyle{fancy}
\fancyhf{}
\fancyhead[L]{\small\sffamily WE-SOLVE \; / \; EXTENDED PROOF EXPOSITION}
\fancyhead[R]{\small\sffamily\reportcode}
\fancyfoot[C]{\small\thepage}
\renewcommand{\headrulewidth}{0.3pt}
\setlength{\emergencystretch}{3em}
\setlength{\parskip}{3pt}
\setlist{itemsep=3pt,topsep=5pt}
\allowdisplaybreaks[2]
\numberwithin{equation}{section}
\setcounter{tocdepth}{1}
\newtheorem{theorem}{Theorem}[section]
\newtheorem{proposition}[theorem]{Proposition}
\newtheorem{lemma}[theorem]{Lemma}
\newtheorem{corollary}[theorem]{Corollary}
\theoremstyle{definition}
\newtheorem{definition}[theorem]{Definition}
\theoremstyle{remark}
\newtheorem{remark}[theorem]{Remark}
\newcommand{\R}{\mathbb R}
\newcommand{\C}{\mathbb C}
\newcommand{\Q}{\mathbb Q}
\newcommand{\N}{\mathbb N}
\newcommand{\Z}{\mathbb Z}
\newcommand{\F}{\mathbb F}
\newcommand{\HH}{\mathbb H}
\newcommand{\eps}{\varepsilon}
\newcommand{\dd}{\,\mathrm d}
\newcommand{\poly}{\operatorname{poly}}
\newcommand{\supp}{\operatorname{supp}}
\newcommand{\rank}{\operatorname{rank}}
\newcommand{\tr}{\operatorname{tr}}
\newcommand{\diag}{\operatorname{diag}}
\newcommand{\diver}{\operatorname{div}}
\newcommand{\pc}{p_{\mathrm c}}
\newcommand{\one}{\mathbf 1}
\newcommand{\sourcecommit}{7ef5c0cc5ca4b0e072b7a050fca09ab24bc6fa3a}
\newcommand{\pin}{\texttt{\sourcecommit}}
\newcommand{\repolink}[2]{\href{https://github.com/openai/math/blob/7ef5c0cc5ca4b0e072b7a050fca09ab24bc6fa3a/#1}{#2}}
\newcommand{\mathsource}[2]{\repolink{preprints/Local-smoothing-for-the-Euclidean-wave-equation-October-25-2026/build/sections/#1.tex}{#2}}
\newcommand{\cssource}[2]{\repolink{preprints/The-Exponential-Time-Hypothesis-is-False-October-25-2026/build/#1.tex}{#2}}
\newcommand{\statusbox}[1]{\begin{quote}\small\noindent\textbf{Attribution and verification boundary.} #1\end{quote}}
\hypersetup{pdfauthor={Source-based exposition prepared with GPT-6 Astra Pro},pdfsubject={We-Solve contribution draft; original research attributed to OpenAI}}
\newcommand{\reportcode}{Math-273}
\title{All-dimensional local smoothing:\\the current, determinant, and propagation arguments\\[4pt]
\large An extended source-pinned proof exposition for Math-273}
\author{Prepared for \texttt{we-solve} from OpenAI's public research artifacts}
\date{26 October 2026 \quad / \quad Expanded edition}
\begin{document}
\section*{All-dimensional local smoothing: the current, determinant, and propagation arguments}
\noindent OpenAI GPT 6 Astra Pro; extended literature report, 26 October 2026.
\section*{Abstract}
This report gives an expanded account of the proof architecture in OpenAI's
October 25, 2026 manuscript on Euclidean local smoothing. It follows the argument
from the symmetric first-order wave system and its analytically continued
positive current, through the quaternionic determinant equation on moving
fibers, to the corrected energy, the boundary and interior barriers, and the
two-circle propagation estimate. The transfer to Euclidean space,
interpolation, and Sobolev summation are then derived explicitly. Separate
sections explain the focusing obstruction, spherical restriction, the Kakeya
maximal and Hausdorff-dimension deductions, Bochner--Riesz means, and Stein's
square function. The central all-dimensional theorem would fully resolve
Math-273; the submission itself remains a literature report giving a detailed
exposition of an externally produced proof, rather than claiming a new discovery
or an independently completed proof audit.

\begin{quote}\textbf{Attribution and verification boundary.} The original analytic construction and research results are attributed
to OpenAI. This exposition is based on the pinned manuscript's TeX sections and
its published formalization documentation. Source lemmas are identified as such;
short deductions are written out here. Reading and explaining the proof is not
the same as independently verifying every analytic identity and limiting
argument. The Lean development was not rebuilt or run in preparing this report.\end{quote}
\paragraph{Reading route.}
Sections~\ref{sec:current}--\ref{sec:propagate} describe the new analytic
mechanism; Section~\ref{sec:transfer} closes the Math-273 deduction.
Sections~\ref{sec:restriction}--\ref{sec:square} explain the principal
consequences. The final sections separate the reported theorem from the
verification performed for this exposition.




\section{The target, the reported theorem, and the scope of this report}
\label{sec:target}
The primary catalogue entry is \emph{Local Smoothing for the Euclidean Wave
Equation}, Math-273, with stable handle
\texttt{prb\_94f7553139864cf78d6311e3db49c212} \cite{catalogue}.
The source snapshot throughout is the \texttt{openai/math} commit
\begin{center}\texttt{7ef5c0cc5ca4b0e072b7a050fca09ab24bc6fa3a}.\end{center}
The original manuscript is dated October 25, 2026. All source-file links in the
bibliography point to this snapshot, not to a moving branch.

Let $n$ denote spatial dimension, put $d=n-1$, and define
\[
 U(T)=e^{-\mathrm iT|D|},\qquad J=(1-\Delta)^{1/2},\qquad
 I=[1,2],\qquad \pc=\frac{2n}{n-1}.
\]
The Sobolev order appearing in the catalogue's supercritical formulation is
\[
 (n-1)\left(\frac12-\frac1p\right)-\frac1p+\eps
       =\frac{n-1}{2}-\frac np+\eps.
\]
Thus its local-smoothing assertion for $p>\pc$ is a time-averaged wave estimate
with this order. The theorem reported in the source includes this range and
also states the critical and subcritical positive-loss bounds.

\begin{theorem}[OpenAI's reported all-dimensional theorem]\label{thm:main}
For every $n\ge2$, $2\le p<\infty$, and $\eps>0$,
\begin{equation}\label{eq:main}
 \|U(T)f\|_{L^p(\R^n\times I)}
 \le C_{n,p,\eps}\|J^{s_+(p)+\eps}f\|_{L^p(\R^n)},\qquad
 s_+(p)=\max\left\{0,\frac{n-1}{2}-\frac np\right\},
\end{equation}
for Schwartz data $f$. At $p=2$, $\eps=0$ is allowed.
\end{theorem}
The manuscript's introduction states this as
\texttt{thm:local-smoothing}; scope document 079 describes the matching
formal target \cite{intro,scope}. The sign in the half-wave is immaterial to
the catalogue comparison: conjugation changes $U(T)$ into $e^{\mathrm iT|D|}$
without changing either norm. For $p>\pc$, $s_+(p)$ is exactly the catalogue's
Sobolev exponent before the positive loss. This is the precise sense in which
the reported theorem is a full solution of Math-273.

This report does not extend the statement to variable-coefficient equations,
curved manifolds, arbitrary mixed space--time exponents, or a lossless endpoint.
Nor does it transfer discovery credit to the preparer. Its purpose is to make
clear what the cited proof does, why its intermediate constructions are needed,
and which analytic implications connect it to the exact catalogue question.

\section{The proof dependency chain and the critical frequency estimate}
\label{sec:roadmap}
The source reduces the theorem to the following annular assertion.
\begin{theorem}[Reported critical annular estimate]\label{thm:annular}
For a smooth compactly supported symbol $\varphi$ away from zero and $N\ge2$,
\begin{equation}\label{eq:annular}
 \|U(T)\varphi(D/N)f\|_{L^{\pc}(\R^n\times I)}
 \le C_{n,\varphi}[\log(2+N)]^{2/\pc}\|f\|_{\pc}.
\end{equation}
\end{theorem}
This is the source theorem labelled \texttt{thm:main} \cite{intro}. It is
stronger than an $N^\eps$ bound, but its logarithmic factor is not asserted to
be optimal for arbitrary data. For example, in dimension three it reads as an
$L^3$ space--time estimate with loss $(\log N)^{2/3}$; in dimension four the
critical exponent is $8/3$.

The hard part is not the final interpolation. It is the production of
\eqref{eq:annular} from an energy living in extra variables. The dependencies
are as follows:
\begin{center}
\begin{tabular}{p{0.28\linewidth}p{0.63\linewidth}}
\toprule
Construction & Mathematical task it performs\\
\midrule
Continued wave current & Retains interference while producing a positive,
column-divergence-free tensor.\\
Rank-one background and fiber equation & Turn the square of a wave into an
energy with the critical $2n/(n-1)$ power.\\
Corrected determinant energy & Remove an indefinite determinant from the
base-variable Laplacian and cancel unwanted cofactor boundary terms.\\
Interior and boundary barriers & Compare energy to the observed wave and
make the remaining positive flux exponentially small.\\
Two-circle comparison & Propagate from observed time back to initial time,
with two logarithmic costs.\\
Transfer and summation & Remove periodic localization and obtain all finite
$p\ge2$ with the stated Sobolev order.\\
\bottomrule
\end{tabular}
\end{center}
The principal notation is $h=N^{-1}$, a small frequency scale, and $a=Kh$,
where $K$ is a large fixed constant. In addition to physical time $T$, the
construction uses an auxiliary two-dimensional time variable $\tau$; its radius
is $s=|\tau|$. Constants are chosen in a definite order, before $h$ tends to
zero. Losing that order could turn an apparent uniform estimate into a
frequency-dependent one.

\section{Why analytic continuation is used instead of packetwise squaring}
\label{sec:filter}
Start with finitely many Fourier modes in a fixed small angular cap of an
annulus, on a fixed spatial torus. The final result is obtained by density,
a finite cap partition, and a period-removal argument. Periodic localization
makes the auxiliary integrations manageable; it is not a new hypothesis on
the Euclidean initial data.

For a mode $e^{\mathrm i(X\cdot\xi-T|\xi|)}$, imaginary translation by
$-\mathrm i\rho(1,\omega)$ gives the multiplier
\begin{equation}\label{eq:lightfilter}
 e^{\rho(\xi\cdot\omega-|\xi|)}
 =\exp\left(-\frac{\rho|\xi|}{2}
                \left|\omega-\frac\xi{|\xi|}\right|^2\right).
\end{equation}
It is a Gaussian directional filter, with angular width approximately
$(N\rho)^{-1/2}$ when $N\rho$ is large. The calculation uses
$|\omega|=|\xi/|\xi||=1$, so
$1-\omega\cdot\xi/|\xi|=|\omega-\xi/|\xi||^2/2$.

The important operation is to filter the \emph{sum} of the modes and then take
its square. Squaring each mode separately first would destroy cross terms.
Those terms are precisely what can make an incoming spherical wave focus.
Consequently a positive quantity built only from independent mode energies
would not automatically control the relevant $L^p$ norm. The source builds a
positive conserved tensor from the continued full wave, rather than replacing
it by such an incoherent sum \cite{current}.

\section{The first-order system and the positive conserved current}
\label{sec:current}
\subsection{Null coordinates expose the appropriate polarization}
Write $X=(x,b)$ with $x\in\R^d$, and introduce
\[
 t=\frac{T+b}{\sqrt2},\qquad s_0=\frac{b-T}{\sqrt2}.
\]
In these coordinates the wave operator is
$2\partial_t\partial_{s_0}+\Delta_x$. A cap mode can be written
\[
 u_{\kappa,v}=c_{\kappa,v}
   e^{\mathrm i\kappa(s_0+v\cdot x-|v|^2t/2)},
 \qquad \kappa>0,
\]
with $v$ confined to a small fixed ball. Let $E_{ab}$ be matrix units in
$n=d+1$ coordinates. The source uses symmetric matrices
\[
 C^t=2E_{00},\qquad C^{x_j}=E_{0j}+E_{j0},\qquad
 C^{s_0}=-\sum_{j=1}^dE_{jj}.
\]
At the mode covector, the negative matrix symbol is
\begin{equation}\label{eq:symbol}
 -C(p)=\kappa
 \begin{pmatrix}|v|^2&-v^{\mathsf T}\\-v&I_d\end{pmatrix}.
\end{equation}
Its quadratic form on $(w_0,w')$ is
$\kappa|w'-vw_0|^2$. It is positive semidefinite and has kernel spanned by
$E(v)=(1,v)$. Thus a vector-valued superposition with polarization $E(v)$
solves the symmetric first-order system. The extra power $\kappa^{d/2}$ in
its normalization is chosen to match the later Gaussian window.

This calculation explains why the source does not simply use the usual
physical stress tensor: the useful positivity is attached to a particular
polarization and an enlarged set of variables. The conservation identities
must then be checked for that enlarged object.

\subsection{Two states and a sum of real squares}
Take two real state vectors $w_1,w_2$ and the quadratic fluxes
\[
 q^\mu=\frac12\sum_{\ell=1}^2 w_\ell^{\mathsf T}C^\mu w_\ell.
\]
The continued wave is evaluated at the associated imaginary displacement.
Combining the states into an auxiliary complex vector $z=(R,\zeta)$ makes
its Gaussian factor proportional to
\[
 \exp\left(-\frac\kappa2|\zeta-vR|^2\right).
\]
The source defines differential fields from the state derivatives and the
quadratic fluxes, and multiplies them by the continued vector wave. In its
notation these are fields $Z_1,Z_2$, and the real tensor is
\begin{equation}\label{eq:realcurrent}
 B=\Re\sum_{\ell=1}^2 Z_\ell\otimes\overline{Z_\ell}.
\end{equation}
This is a sum of four real rank-one squares, hence positive semidefinite.

Positivity alone is not conservation. There are two further identities. The
first-order equation gives $\diver Z_\ell=0$. The acceleration contribution
in $\sum_\ell\nabla_{Z_\ell}\overline{Z_\ell}$ is purely imaginary.
Taking the real part therefore leaves every real column of $B$ divergence
free. Both identities depend on the specific paired-state construction;
they are not generic properties of analytically continued functions.
The full differentiated formulas are in the source current section
\cite{current}.

There is also a rank obstruction worth making explicit. Although $B$ is
positive, it is a sum of only four real squares. In larger ambient dimensions
its ordinary determinant vanishes. The proof therefore cannot end by applying
an off-the-shelf estimate directly to $\det B$. Instead it uses a component of
this tensor as the density in a different determinant equation.

\section{The lifted coordinates, tangency, and wave detection}
\label{sec:lift}
The next transformation has to preserve the conservation identities while
making the directional Gaussian explicit. With a passive auxiliary variable
$\eta$, the source introduces the quadratic holomorphic flux map
\[
 \Theta(z,\eta)_\mu=\frac12z^{\mathsf T}C^\mu z+c_\mu\eta.
\]
For the wave system its null-coordinate version is
\[
 Y=R^2,\qquad P=R\zeta,\qquad S_1=\eta-\frac12\zeta\cdot\zeta.
\]
Consequently
\begin{equation}\label{eq:firstintegral}
 W=P/Y,\qquad
 \Phi=S_1+\frac{P\cdot P}{2Y}=\eta,
 \qquad y=|Y|.
\end{equation}
The quadratic products here are the auxiliary complex bilinear products,
not squared absolute values. Confusing these two products would destroy the
holomorphic cancellation.

Under a nonlinear change of coordinates, divergence-free tensor columns do
\emph{not} in general remain divergence free merely by applying a Jacobian.
The source verifies the transformed conservation law using holomorphicity and
the equality of the two state blocks. The real Jacobian factor and the
pushforward tensor density are part of this verification. This is a substantive
step in the construction, not a coordinate convention that can be omitted.

The resulting wave tensor can be expressed as a positive quaternionic Hermitian
rank-one tensor $A_0=\chi\chi^*$. The vector $\chi$ sums the full wave against
windows of the form
\begin{equation}\label{eq:window}
 m_{\kappa,v}(y,W)
   =\left(\frac\kappa y\right)^{d/2}
       e^{-\kappa y|W-v|^2/2}.
\end{equation}
The accompanying polarization, in the source's null coordinates, is
\[
 V^v=\left(1,\frac{v+W}{2},-\frac{W\cdot v}{2}\right).
\]
A direct substitution gives $\partial\Phi\,V^v=0$. Hence the tensor is
tangent to the level sets of $\Phi$, a fact later used both for cutoffs and
for the rank-one background.

The physical time component is still sensitive to the original wave. More
precisely, the source writes $(A_0)_{TT}$ as the sum of squares of two scalar
windowed waves. In the small cap and near the central window, the first scalar
window is invertible on the cap. The corresponding inverse multiplier estimate
is what permits a lower bound by the original wave after suitable averaging.
No claim is being made that a single arbitrary Gaussian sample recovers the
wave with a universal constant. The location of the window and the scale
$y\asymp s$ matter \cite{current,barrier}.

\section{The quaternionic Hessian calculus behind the energy}
\label{sec:calculus}
\subsection{Why the determinant is an auxiliary spatial determinant}
For each quaternionic coordinate the source uses three real coordinate slots,
with quaternionic units $e_0=1,e_1=\mathrm i,e_2=\mathrm j$. Its gradient and
Hessian operators have the schematic, and in its conventions literal, forms
\[
 G_a(u)=\sum_{\alpha=0}^2e_\alpha\partial_{a\alpha}u,
 \qquad
 H_{ab}(u)=\sum_{\alpha,\beta=0}^2
 e_\alpha\overline{e_\beta}\,
             \partial_{a\alpha}\partial_{b\beta}u.
\]
The matrix $H(u)$ is quaternionic Hermitian and $\tr H(u)=\Delta u$.
Positivity of a real Hessian implies positivity of this Hermitian Hessian.
The product rule is
\begin{equation}\label{eq:producthess}
 H(uv)=uH(v)+vH(u)+G(u)G(v)^*+G(v)G(u)^*.
\end{equation}
In the appropriate auxiliary holomorphic variables,
$H(|F|^2)$ is a positive Jacobian square. In particular, when $F$ is the
single scalar first integral $\Phi$, the background Hessian has rank one.

The determinant used below is the Moore determinant of a Hermitian matrix.
On positive matrices it is the product of the real eigenvalues. Its cofactor,
written $\operatorname{cof}(A)$, satisfies
\[
 \operatorname{cof}(A)=\det(A)A^{-1}
\]
when $A$ is invertible, and extends polynomially to singular matrices. For
positive semidefinite $A$ the cofactor is positive semidefinite. At a singular
matrix, its image is confined to null directions of $A$.
These facts allow the proof to exploit degeneracy rather than erase it by
assuming every matrix is invertible \cite{calculus}.

\subsection{The cofactor divergence identity}
A central identity is that the real columns of $\operatorname{cof}(H(u))$ are
divergence free. This is the analogue of the divergence structure of Hessian
cofactors in real Monge--Ampere calculus. The source proves it through
polarization, including at singular and indefinite Hessians.

The useful algebraic observation is that the cofactor is homogeneous of degree
one less than the matrix size. Polarize it on rank-one symbols $g_jg_j^*$.
The resulting mixed cofactor annihilates the relevant span of the $g_j$.
Testing on exponential functions turns differentiation into multiplication by
the corresponding frequencies; the annihilation gives the cancellation.
Since the desired statement is a differential polynomial identity, the
polarized calculation proves the identity for general smooth $u$.

It follows that, for a tensor with divergence-free real columns, pairing it
with $H(u)$ is a real divergence. This permits the integrations by parts that
produce the corrected energy. The identities are not restricted to a formal
interior calculation: the boundary contractions still have to be evaluated.

\subsection{Schur complement and normal compression}
For a Hermitian block matrix with positive spatial block $D$,
\[
 H=\begin{pmatrix}\ell&c^*\\c&D\end{pmatrix},\qquad
 \sigma=\ell-c^*D^{-1}c,
\]
one has $\det H=\det D\,\sigma$ and
$H-\sigma e_Te_T^*\succeq0$. This isolates the possibly negative scalar
part of the full space--time Hessian.

A second identity identifies a normal cofactor contraction with the determinant
of the compression to the normal's orthogonal complement. It is this
compression, not the full determinant by itself, that controls a cofactor
boundary flux. A common tangential null direction can make that compression
singular and the flux vanish. The fiber geometry is designed to supply exactly
such a direction.

\section{The moving fibers and the Dirichlet equation}
\label{sec:dirichlet}
\subsection{The geometry is part of the proof}
Let $z=(T,\tau)$ denote the base variables and $s=|\tau|$. For a small fixed
$w_0$, the source's defining function is
\begin{equation}\label{eq:fiber}
 r=\frac{|W|^2}{w_0^2}+100|Q|^2-1,
 \qquad Q=\Phi/a,
 \qquad \Omega_z=\{r<0\},\quad a=Kh.
\end{equation}
The real spatial directions are periodic. On the relevant annulus in $s$,
the change of variables gives $y\asymp s$, an auxiliary-volume factor
comparable to $a^2s^{2d}$, and spatial Hessian scales
\[
 H_{\mathrm{fib}}(r)\succeq c s^{-2}I,
 \qquad
 \det H_{\mathrm{fib}}(r)\asymp a^{-2}s^{-2d}.
\]
The complete Hessian has a null direction which is tangent to the fiber wall
and also belongs to the kernel of $H(|\Phi|^2)$. This is the shared direction
needed in the normal-compression argument.

The positive tensor is regularized by spatial convolution and a cutoff in
$Q$, followed by a small cofactor term. In the source's normalization it has
the structure
\[
 A=A_w+h^{50n}\operatorname{cof}(H(r)),\qquad
 A_w=g(Q)(J_h*_X A_0),\qquad \mu=A_{TT}>0.
\]
Tangency to $\Phi$ makes the cutoff compatible with conservation. The cofactor
term preserves the divergence law and has no outward wall flux of the type
that must be controlled later. These choices prevent regularization from
introducing an uncontrolled error in the central identity \cite{energy}.

\subsection{A rank-one background creates the critical exponent}
Set $B_0=\lambda|\Phi|^2$, with $\lambda=h^{-30n}$ in the source. Its spatial
Hessian $B_{\mathrm{fib}}$ has one nonzero eigenvalue $\beta\asymp\lambda$.
The construction solves
\begin{equation}\label{eq:MA}
 D_u=H_{\mathrm{fib}}(B_0+u)>0,
 \qquad \det D_u=\lambda\mu,
 \qquad u|_{\partial\Omega_z}=0.
\end{equation}
The solution satisfies $u\le0$. Tangency gives $A:H(B_0)=0$, so this background
has no cost in the current pairing, even though it contributes curvature to
the spatial determinant.

The exponent can already be predicted. Adding isotropic curvature $\alpha I$
to the rank-one background changes the determinant to
\begin{equation}\label{eq:rankone}
 \det(B_{\mathrm{fib}}+\alpha I)=(\beta+\alpha)\alpha^{n-1}.
\end{equation}
In the controlled range, the $\lambda$ factor cancels the right side of
\eqref{eq:MA}; the necessary curvature is of order $\mu^{1/(n-1)}$.
The energy pairs a depth of this order with $\mu$, creating
$\mu^{n/(n-1)}$. Since $\mu$ detects a squared wave, the resulting wave
power is $2n/(n-1)=\pc$. An unstructured full-rank background would not give
this same determinant scaling.

\subsection{What the existence theorem supplies}
The source proves a Dirichlet theorem for periodic-real fibers with a smooth
strictly Hessian-convex defining function, a semipositive background, and a
positive density above its background determinant. It supplies a unique smooth
admissible solution, comparison barriers of the form
\[
 b r\le u\le b' r,\qquad b\ge b'>0,
\]
and smooth dependence on locally varying parameters \cite{dirichlet}.
These are substantial analytic inputs. They are not assumed merely because
one can write the equation.

The proof uses a continuity path between a known barrier density and the target
density. Linearizing the logarithmic determinant at a positive matrix $D$
gives the operator $D^{-1}:H(w)$, whose real principal symbol is positive.
The logarithmic determinant is concave. Comparison follows by averaging the
linearizations between two admissible potentials and applying the maximum
principle.

For the a priori estimates, the defining function gives a strict subsolution
and controls the boundary gradient. Tangential second derivatives at the
boundary are expressed through $u=\rho r$. Mixed normal--tangential derivatives
are bounded using tangential rigid vector fields and barriers involving
$|x|^2$ and $-r$. The restriction of the Hessian to the tangent complement has
a uniform positive lower bound. The determinant equation and Schur complement
then recover the missing pure normal derivative estimate.

The source temporarily introduces an additional periodic real slot to permit
the full quaternionic rotation argument. Uniqueness and invariance under
translations in that slot imply that the solution descends to the original
three-slot setting. Interior second-derivative control, concavity, and uniform
ellipticity permit regularity estimates and bootstrapping. Pulling moving
fibers back to a fixed domain, followed by the implicit-function argument,
gives smooth parameter dependence.

There is an important uniformity distinction. This theorem provides smooth
solutions for each admissible regularized problem. It does not by itself give
the frequency-uniform wave estimate. The estimates with controlled dependence
on $h$, $s$, and $a$ are proved by the explicit barriers below.

\section{Why the naive energy must be corrected}
\label{sec:energy}
The natural positive quantity is
\[
 X_0(z)=\int_{\Omega_z}(-u)\mu.
\]
A calculation with the full Hessian of $B_0+u$ and its spatial block $D_u$
produces a Schur decomposition
\[
 H(B_0+u)=\sigma e_Te_T^*+\Sigma,\qquad \Sigma\succeq0.
\]
Since $\det D_u=\lambda\mu$ and $A:H(B_0)=0$, its pairing contains
\begin{equation}\label{eq:unwanted}
 A:H(u)=A:\Sigma+\lambda^{-1}\det H(B_0+u).
\end{equation}
The first term has a sign; the full determinant need not. Dropping the last
term would therefore be invalid. The proof adds an energy whose variation
cancels it exactly.

Define the fiber determinant energy
\[
 E_z(v)=\lambda^{-1}\int_{\Omega_z}v
       \int_0^1\det H_{\mathrm{fib}}(B_0+\theta v)\dd\theta,
 \qquad \mathcal F=X_0+E_z(u).
\]
The interpolation in $\theta$ is the essential part of the definition. It is
not an arbitrary correction chosen only to make the final inequality true.
The cofactor differentiation formula, combined with its divergence identity,
turns the variation of this integral into the determinant term occurring in
\eqref{eq:unwanted} \cite{energy,calculus}.

\subsection{The cancellation on a moving domain}
The integrations by parts are carried out over the total moving-fiber domain,
using a test function of the base variable $z$. Treating the fiber as fixed
while differentiating in $z$ would miss boundary contributions. Writing
$\rho=u/r$ at the wall, the resulting identity has the form
\begin{equation}\label{eq:energyidentity}
 \Delta_z\mathcal F+\int_{\Omega_z}A:\Sigma
 =\int \rho\,G(r)^*A_wG(r)\,\delta(r)\dd\zeta.
\end{equation}
Here the right-hand side is the boundary integral written by the coarea
convention with $\delta(r)$. It is nonnegative.

Why are there not additional cofactor terms? At the wall, the Hessian appearing
in the relevant normal compression is built from the rank-one background and
a multiple of $H(r)$. Their common tangential null direction makes the
compressed determinant vanish. The cofactor normal--normal contraction is
therefore zero. The regularizing cofactor current has the same zero-flux
property. What remains is the physical wave flux $A_w$ on the right of
\eqref{eq:energyidentity}. This precise boundary cancellation is one of the
main reasons for the elaborate fiber construction.

\subsection{The corrected energy remains positive}
The source derives the identity
\begin{equation}\label{eq:energycomparisonidentity}
 \mathcal F=\frac{n-1}{n+1}X_0+
 \frac{1}{\lambda n(n+1)}
 \int_{\Omega_z}G_{\mathrm{fib}}(u)^*
              \operatorname{cof}(D_u)G_{\mathrm{fib}}(u),
\end{equation}
and the two-sided comparison
\begin{equation}\label{eq:energycomparison}
 \frac{n-1}{n+1}X_0\le\mathcal F\le\frac n{n+1}X_0.
\end{equation}
Thus the cancellation does not destroy positivity or change the relevant
size of the energy.

The rank-one hypothesis makes this computation particularly transparent. Along
the interpolation segment,
\[
 \det((1-\theta)B_{\mathrm{fib}}+\theta D_u)
 =\theta^n\det D_u+
  (1-\theta)\theta^{n-1}\operatorname{cof}(D_u):B_{\mathrm{fib}}.
\]
Terms with two or more columns from the rank-one background vanish. Integrating
these two powers of $\theta$ and then using the cofactor integration by parts
gives \eqref{eq:energycomparisonidentity}; positivity and the companion estimate
give \eqref{eq:energycomparison}. This is the point where the determinant
structure yields an actual usable energy rather than just a formal identity.

\section{Boundary decay and interior wave recovery}
\label{sec:barriers}
The next objective is to replace the right side of
\eqref{eq:energyidentity} by a small multiple of $\mathcal F$ itself.
An upper bound for the boundary flux alone is not enough: one also needs a
lower bound for the energy that detects the same wave.

\subsection{A barrier that controls the boundary normal derivative}
The $Q$ cutoff is supported near the center of its allowed disk. At the portion
of the wall where that cutoff is nonzero, \eqref{eq:fiber} forces $|W|$ to be
bounded away from zero. The frequency directions satisfy $|v|\le\gamma$ for
a much smaller fixed $\gamma$. Hence $|W-v|$ is uniformly bounded below there,
and \eqref{eq:window} gives exponential decay $e^{-c_0s/h}$.

To make a comparison function robust under differentiation, the source uses a
positive spatially smoothed density
\[
 P_*(X)=\bigl(J_{Ds}*|g_T|^2(X)\bigr)^{1/d},
 \qquad g_T=U(T)f,
\]
with a broad smoothing kernel and a large fixed $D$. Its first two derivatives
are controlled relative to itself at scale $Ds$, and
\begin{equation}\label{eq:psmooth}
 \int P_*^{d+1}\lesssim\|g_T\|_{\pc}^{\pc}.
\end{equation}
The broad scale makes the mixed derivative errors in $H_{\mathrm{fib}}(rP_*)$
small compared with its positive principal part. A sufficiently large
multiple of $(s/h)^m rP_*$ is therefore a determinant subsolution, where $m$
is a fixed exponent depending on the construction.

Comparison with $u$ bounds the boundary ratio $\rho=u/r$ and yields both an
upper bound for the energy and an exponentially small flux estimate. With
$\sigma=s/h$, the estimates have the schematic quantitative form
\begin{align}
 \mathcal F(T,\tau)&\le C a^2\sigma^m\|g_T\|_{\pc}^{\pc},
                                                    \label{eq:upperF}\\
 \Delta_z\mathcal F(T,\tau)&\le
 C(a/s)^2\sigma^m e^{-c_0\sigma}\|g_T\|_{\pc}^{\pc}.
                                                    \label{eq:fluxF}
\end{align}
The symbols $C,m,c_0$ here stand for the fixed constants established by the
source's barrier estimates, not quantities allowed to grow with frequency
\cite{barrier}.

\subsection{An interior comparison recovers the critical norm}
Near the central window, the inverse scalar window multiplier gives a lower
bound for the regularized density on small interior balls. In the source's
scale notation, a useful form of the lower density is
\[
 m_X\gtrsim (hs)^{-d}\sigma^{-m}e^{-C\gamma^2\sigma}|g_T(X)|^2.
\]
The ball radius is comparable with $h$, and it lies inside the cutoff-one
region when $K$ is sufficiently large.

On such a ball, use a negative quadratic potential with curvature proportional
to $m_X^{1/d}$. Equation \eqref{eq:rankone}, together with the density ceiling
in the construction, makes its determinant no larger than $\lambda m_X$.
The comparison principle forces the depth of $u$ at the center to be at least
a constant multiple of $h^2m_X^{1/d}$. Multiplication by the density and
integration over the interior windows give
\begin{equation}\label{eq:lowerF}
 \mathcal F(T,\tau)\ge
 c a^2\sigma^{-m}e^{-C_1\gamma^2\sigma}\|g_T\|_{\pc}^{\pc}.
\end{equation}
The power $\pc$ is the rank-one determinant mechanism of
Section~\ref{sec:dirichlet}; the factor $a^2$ comes from the fiber-volume
normalization and the wave-window scaling.

Combining \eqref{eq:fluxF} with \eqref{eq:lowerF}, first choose $\gamma$ so small
that $C_1\gamma^2<c_0/2$. Then choose $K$ large enough that the remaining
polynomial factors in $\sigma$ are absorbed by exponential decay. This gives
\begin{equation}\label{eq:scalarinequality}
 \mathcal F\ge0,\qquad
 \Delta_{T,\tau}\mathcal F
 \le C s^{-2}e^{-cs/h}\mathcal F,
 \qquad 2a\le s\le3.
\end{equation}
At the two particular radii needed later, the same estimates give
\begin{align}
 \mathcal F(T,\tau)&\ge c a^2\|U(T)f\|_{\pc}^{\pc}
                         &&(|\tau|=2a),\label{eq:innerobserve}\\
 \mathcal F(0,\tau)&\le C a^2\|f\|_{\pc}^{\pc}
                         &&(|\tau|=4a).\label{eq:outerinitial}
\end{align}
The constants may depend on the fixed cap and $K$, but not on $h$. The order
of choices is essential: shrinking the cap after choosing frequency-dependent
constants would not prove the desired uniform estimate.

\section{The two-circle propagation argument and its two logarithms}
\label{sec:propagate}
Let $H(T,s)$ be the circular average of $\mathcal F(T,\tau)$ over $|\tau|=s$.
The analytic construction has now been reduced to a scalar comparison problem
in one physical time and two auxiliary time coordinates.

The source proves a general propagation lemma for a nonnegative smooth
function $E$ satisfying $\Delta E\le V(s)E$ on a cylindrical annulus
$b\le s\le3$. If
\[
 \int_b^3 s\log(s/b)V(s)\dd s\le M,
 \qquad \sup_{1/2\le s\le3/2}V(s)\le M,
\]
then its circular mean obeys
\begin{equation}\label{eq:propagation}
 H(0,r)\ge c(M)
 \frac{\log(r/b)}{[\log(3/b)]^2}
             \int_1^2 H(T,b)\dd T,
 \qquad b\le r\le1.
\end{equation}
The detailed scalar lemma and its boundary comparisons are in
\cite{propagation}. The two logarithms arise from two distinct radial passages,
not from an unexplained final loss.

\subsection{The radial equation becomes a one-dimensional problem}
The relevant radial comparison equation is
\[
 u''(s)+s^{-1}u'(s)=(V(s)+\kappa^2)u(s).
\]
Set $x=\log(s/b)$. Then the left side becomes $s^{-2}u_{xx}$, so the
comparison is a second-order equation on a logarithmic interval. A Volterra
integral equation and a Gronwall estimate control its normalized positive
solutions by the integral condition on $V$. When $V=\kappa=0$, the radial
solutions are affine in $\log s$; this already explains why logarithms, rather
than a power of $b$, govern communication between concentric circles.

\subsection{Outward propagation, time transfer, and return}
For the outward passage, integrate the circular mean against a positive
sine weight in physical time. Integration by parts gives an inequality with
a fixed spectral term $\kappa^2$; the boundary time contributions have the
favorable sign. The radial comparison moves observed mass from radius $b$
to a fixed-radius annulus, with a factor comparable to
$1/\log(3/b)$.

At fixed radius the potential is bounded. A local positive comparison or
Harnack argument moves the resulting information in physical time to a
neighborhood of initial time. This step has a fixed-scale constant, not a
new small-radius logarithm.

The return passage uses a second positive time weight and a radial function
vanishing at the inner boundary. Moving back from unit radius to $r$ costs
another factor $\log(r/b)/\log(3/b)$. Multiplying the two radial costs gives
\eqref{eq:propagation}. In particular, when $r=2b$, the numerator is only
$\log2$, so both inverse logarithms remain.

\subsection{Closing the cap estimate}
Apply the lemma with $b=2a$ and
$V(s)=Cs^{-2}e^{-cs/h}$. Its weighted integral is bounded uniformly as
$h\downarrow0$: after $s=bt$, the exponential is $e^{-2cKt}$, so the integral
is controlled by an integrable function of $t$ independent of $h$.
Taking $r=4a=2b$, and using \eqref{eq:innerobserve} and
\eqref{eq:outerinitial}, yields
\[
 a^2\int_1^2\|U(T)f\|_{\pc}^{\pc}\dd T
 \lesssim [\log(2/h)]^2 a^2\|f\|_{\pc}^{\pc}.
\]
The $a^2$ factors cancel. Taking a $\pc$-th root gives the periodic fixed-cap
version of \eqref{eq:annular}. The source observes that two logarithms can
occur for the scalar comparison problem itself. This does not prove that the
wave estimate's logarithmic power is optimal: the wave-generated energy is
only one special class of scalar functions satisfying the inequality.

\section{Removing the torus and deriving all Sobolev estimates}
\label{sec:transfer}
\subsection{A finite angular cover}
The cap aperture is fixed independently of $h$. A finite smooth partition
of the frequency annulus therefore reduces general annular data to the cap
estimate without introducing a growing number of caps. Each cap cutoff has a
uniformly integrable spatial kernel after rescaling. Summing the finitely many
bounds preserves the logarithmic power.

\subsection{The off-cone kernel controls periodization errors}
The Euclidean annular wave kernel is
\[
 K_h(T,X)=(2\pi h)^{-n}\int
 e^{\mathrm i(X\cdot\eta-T|\eta|)/h}\varphi(\eta)\dd\eta.
\]
For $T\in[1,2]$ and $|X|\ge3$, the gradient of the unscaled phase has magnitude
at least $|X|-2\ge|X|/3$. Repeated integration by parts therefore gives, for
every sufficiently large integer $M$,
\begin{equation}\label{eq:offcone}
 |K_h(T,X)|\le C_Mh^M(1+|X|)^{-M}.
\end{equation}
The amplitude is supported away from zero, so its differentiated phase and
symbol terms are uniformly controlled.

Tile Euclidean space by unit cubes $Q$. For points in $Q$, split the input
into its restriction to $10Q$ and its complement. The far contribution is
bounded using \eqref{eq:offcone} and Young's inequality. Periodize the near
input on a torus of a fixed sufficiently large side, for example the source's
side $100$. All nonzero translates of the kernel are off cone on $Q$ and
contribute an $O(h^M)$ error. Apply the periodic estimate to the near part,
raise to the $\pc$-th power, and sum over cubes. The enlarged cubes have
bounded overlap. This proves the Euclidean annular estimate without letting
the torus size depend on frequency \cite{transfer}.

\subsection{Interpolation gives exactly the stated derivative threshold}
At $p=2$, Plancherel and the length-one observation interval give a uniform
bound. At $p=\infty$, the annular wave kernel has $L^1$ norm $O(N^{d/2})$.
One way to see the latter is to split the annulus into $O(N^{d/2})$ angular
caps of aperture $N^{-1/2}$; each rescaled cap kernel has uniformly bounded
$L^1$ norm after its linear phase is removed.

For $2\le p\le\pc$, interpolate the $L^2$ estimate with
\eqref{eq:annular}. If $1/p=(1-\theta)/2+\theta/\pc$, then
$2\theta/\pc=d(1-2/p)$. No positive power of $N$ appears.
For $\pc\le p<\infty$, interpolate the critical estimate with $L^\infty$.
The weight of the critical estimate is $\pc/p$, so the power of $N$ is
\[
 \frac d2\left(1-\frac{\pc}{p}\right)=\frac d2-\frac np.
\]
Together these give
\begin{equation}\label{eq:allp}
 \|U(T)P_Nf\|_{L^p_{T,X}}
 \lesssim N^{s_+(p)}[\log(2+N)]^{\beta(p)}\|f\|_p,
 \quad
 \beta(p)=\min\{d(1-2/p),2/p\}.
\end{equation}
The two formulas for $\beta$ agree at $\pc$.

\subsection{The positive loss is what makes frequency summation immediate}
Set $g=J^{s_+(p)+\eps}f$ and take a smooth dyadic frequency partition.
On the annulus $N=2^j$, a rescaled multiplier estimate gives
\[
 \|U(T)P_{2^j}f\|_{L^p_{T,X}}
 \lesssim 2^{-j\eps}(1+j)^{\beta(p)}\|g\|_p.
\]
The series of coefficients converges for every $\eps>0$. The remaining
low-frequency multiplier is uniformly bounded by standard smooth multiplier
estimates. Summing gives \eqref{eq:main}. At $p=2$, use the unitary estimate
directly, avoiding a needless loss.

The source also gives a finer logarithmic Sobolev summation. Put
$\delta_p=1/2-1/p$ and $B_*(p)=\beta(p)+\delta_p$. A bound with the exact
power $J^{s_+(p)}$ and additional weights
$[\log(2+J)]^B[\log(2+\log(2+J))]^C$ follows when $B>B_*(p)$, or when
$B=B_*(p)$ and $C>\delta_p$ \cite{transfer}. The reason is that a spatial
Littlewood--Paley square sum costs an $\ell^2$ norm across scales; Holder from
$\ell^p$ requires the remaining weight in $\ell^r$, $1/r=\delta_p$.
At the boundary one obtains the convergence test
$\sum j^{-1}(\log j)^{-Cr}<\infty$. These summation conditions are not
claimed to be optimal for every $p$.

\section{The focusing test and what the endpoint still does not say}
\label{sec:sharpness}
The source supplies a radial test that explains both the critical exponent
and the unavoidable logarithm \cite{sharpness}. Fix $t_0=3/2$ and take
\[
 \widehat f_N(\xi)=e^{\mathrm it_0|\xi|}\psi(|\xi|/N),
\]
where $\psi$ is a nonzero nonnegative smooth bump in a compact subinterval of
$(1,2)$. The data are an incoming spherical wave chosen to focus at $t_0$.
Plancherel gives $\|f_N\|_2\asymp N^{n/2}$, while stationary phase on the
sphere gives $\|f_N\|_\infty\lesssim N^{(n+1)/2}$. Interpolation yields
\[
 \|f_N\|_p\lesssim N^{(n+1)/2-1/p}.
\]

Write $T=t_0+u$. On an outgoing shell $|r-u|\lesssim N^{-1}$ with
$A/N\le u\le1/4$, stationary phase gives wave magnitude comparable from
below to $N^{(n+1)/2}u^{-d/2}$, after choosing the shell so the principal
oscillatory term dominates its opposite phase and remainder. The shell volume
is comparable to $N^{-1}u^d$. Therefore
\begin{equation}\label{eq:focusintegral}
 \frac{\|U(T)f_N\|_{L^p_{T,X}}^p}{\|f_N\|_p^p}
 \gtrsim \int_{A/N}^{1/4}u^{d-dp/2}\dd u.
\end{equation}
At $p=\pc$ the integrand is $u^{-1}$, giving a lower logarithmic loss
$(\log N)^{1/\pc}$. Above $\pc$, the lower power is
$N^{d/2-n/p}$. Below $\pc$, this radial test gives a constant lower bound.

Thus the critical unrestricted annular logarithmic exponent lies between the
reported lower value $1/\pc$ and upper value $2/\pc$. The source has further
weak-norm and angular-regularity refinements; they are not needed for the
Math-273 implication, and no stronger endpoint conclusion is inferred here.
In particular, taking $\eps\downarrow0$ in \eqref{eq:main} is not legitimate
without uniform control of $C_{n,p,\eps}$.

\section{From a focusing shell to spherical extension}
\label{sec:restriction}
Let
\[
 Eg(X)=\int_{S^{n-1}}e^{\mathrm iX\cdot\omega}g(\omega)\dd\sigma(\omega)
\]
be the spherical extension operator. The source derives a local critical
extension estimate by using a modulated time average of annular waves
\cite{kakeya-proof}. This is a useful bridge: spherical data emerge from a
wave whose incoming radial phase is chosen to focus.

Take an operator of the form
\[
 M_N=\int\psi(T)e^{\mathrm iNT}U(T)\chi(|D|/N)\dd T.
\]
The annular local-smoothing estimate bounds its $L^{\pc}$ norm by a power of
$\log N$. Its kernel on a fixed radial shell has leading stationary-phase
term $N^{d/2}|X|^{-d/2}e^{\mathrm iN|X|}$ times a smooth amplitude.
Apply this to incoming data
$f_N(r\omega)=e^{-\mathrm iNr}\psi(r)g(\omega)$ and evaluate near the focus
at scale $X/N$. The oscillating radial factors cancel, and the rescaled
output tends to a nonzero multiple of $Eg(-X)$.

On a ball of radius $R$, the source controls the error by a polynomial in
$R$ divided by $N$ and chooses $N$ as a sufficiently large polynomial of $R$.
The scaling factors cancel at the critical exponent because
$n/\pc=d/2$. This produces
\begin{equation}\label{eq:localextension}
 \|Eg\|_{L^{\pc}(B_R)}\lesssim
 [\log(2+R)]^{A}\|g\|_{L^{\pc}(S^{n-1})},\qquad A=2/\pc.
\end{equation}
Modulating $g$ translates the spatial ball, so the bound is uniform in its
center. That uniformity is important in the later tube argument.

The source then applies an epsilon-removal theorem to obtain global diagonal
extension for every exponent strictly larger than $\pc$; this is an external
analytic implication explicitly used in the source, not a new elementary
consequence proved solely by summing balls \cite{global-extension}.
Interpolating the diagonal bounds with $L^1(S^{n-1})\to L^\infty(\R^n)$
gives the mixed strict range
\begin{equation}\label{eq:restrictionrange}
 \|Eg\|_{L^b(\R^n)}\lesssim\|g\|_{L^a(S^{n-1})},\qquad
 b>\frac{2n}{n-1},\quad
 b>\frac{n+1}{n-1}a',
\end{equation}
for $1<a\le\infty$ and finite $b$.

For clarity, the exponent calculation when $a<b$ can be made explicit.
Choose a diagonal exponent $s=1+b/a'$. Interpolation with weight $\theta=s/b$
gives $1/b=\theta/s$ and $1/a=1-\theta+\theta/s$.
The two strict inequalities ensure that a valid diagonal exponent can be
used, with a small adjustment if needed. If $a\ge b$, the finite surface
measure gives the necessary embedding. The source also obtains the
corresponding truncated-cone bounds by slicing in its radial parameter.

The inspected Math-96 entry includes equality on one of its exponent boundaries.
The released statement \eqref{eq:restrictionrange} has strict inequalities.
Accordingly this exposition does not mark the entire Math-96 formulation
resolved from that statement alone. Taking a boundary limit of constants which
may diverge is not a proof of the boundary estimate \cite{restriction-catalogue}.

\section{Tube overlap, the Kakeya maximal operator, and Hausdorff dimension}
\label{sec:kakeya}
\subsection{Random signs extract tube overlap from extension}
Consider unit tubes $T_j$ of radius comparable to $\delta$ with
$\delta$-separated directions and nonnegative weights $a_j$.
Let $q=\pc/2=n/(n-1)$. The source obtains
\begin{equation}\label{eq:tubes}
 \left\|\sum_j a_j\one_{T_j}\right\|_q
 \lesssim[\log(2/\delta)]^{2A}
          \left(\delta^{n-1}\sum_j a_j^q\right)^{1/q}.
\end{equation}
To understand this implication, place disjoint small spherical caps around
the tube directions. Modulate each cap by a phase determined by its tube
center. On the appropriate rescaled tube, the cap extension has magnitude
bounded below by a constant times $\delta^{n-1}$.

Apply \eqref{eq:localextension} to a random-sign sum of these modulated caps,
with coefficient $\sqrt{a_j}$. Averaging over the signs retains a square
function; its $\pc/2$ power dominates the weighted tube-overlap expression.
The spatial dilation is $\delta^{-2}$, and the critical scaling again cancels:
$d\pc-2n=0$. Grouping tube centers into fixed-size spatial cells allows the
uniform local extension estimate to be applied without an unbounded-radius
constant. The logarithmic power doubles because the tube overlap is obtained
from a square of the extension. This gives \eqref{eq:tubes}
\cite{kakeya-proof}.

\subsection{Duality gives the maximal bound}
Define the Kakeya maximal operator by taking the largest average of $|f|$
over unit $\delta$-tubes in each direction. Discretize the sphere into
$\delta$-separated directions, apply duality to \eqref{eq:tubes}, and then
pass to the spherical measure formulation. Since $q'=n$, the result is
\begin{equation}\label{eq:kmax}
 \|K_\delta f\|_{L^n(S^{n-1})}
 \lesssim [\log(2/\delta)]^{2A}\|f\|_{L^n(\R^n)}.
\end{equation}
In particular, every fixed positive power loss $\delta^{-\eps}$ is allowed.
The source's formal scope states this latter epsilon-loss maximal theorem.
One should not automatically identify every sharper logarithmic estimate in
the manuscript with the exact statement checked by the published comparator.

Applying \eqref{eq:kmax} to an appropriate neighborhood of a Kakeya set gives
a volume lower bound and full Minkowski dimension. But full Minkowski dimension
does not by itself prove full Hausdorff dimension. The proof needs a genuinely
multiscale covering argument.

\subsection{The multiscale step cannot be replaced by a neighborhood bound}
Let a set contain a unit segment in every direction, and cover it by small
balls. Partition the balls according to their radii, for example
$e^{-2^{k+1}}<r_j\le e^{-2^k}$. Let $G_k$ be the union of their doubled balls,
$\delta_k=e^{-2^{k+1}}$, and $v_k$ the sum of their $n$-th powers of radii.
The total length of a covered segment forces a pointwise directional inequality
of the form
\[
 1\lesssim\sum_k K_{\delta_k}\one_{G_k}.
\]
After integrating over directions and using \eqref{eq:kmax},
\[
 1\lesssim\sum_k 2^{2Ak}v_k^{1/n}.
\]
Holder's inequality converts this to a lower bound for the cost of the cover
in the gauge
\begin{equation}\label{eq:gauge}
 h_\gamma(r)=r^n[\log(1/r)]^{2An}
                  [\log\log(1/r)]^\gamma,
 \qquad \gamma>n-1.
\end{equation}
The second logarithmic weight is selected so that the complementary tail
$\sum_k k^{-\gamma/(n-1)}$ converges. As the maximum covering radius tends
to zero, its tail tends to zero, forcing the gauge cost to diverge. Since for
every $s<n$ the gauge in \eqref{eq:gauge} is eventually smaller than $r^s$,
this implies Hausdorff dimension at least $s$ for all $s<n$, hence exactly $n$.
The argument uses covers, so it does not assume measurability of the original
set.

This supplies the full-dimensional Hausdorff conclusion relevant to Math-95
for every $n\ge4$, as well as the lower dimensions included in the source.
It says full dimension, not positive Lebesgue measure
\cite{kakeya-proof,kakeya-catalogue}.

\section{Bochner--Riesz means: the shell argument}
\label{sec:br}
For the multiplier
\[
 B_r^\alpha=(1-|D|^2/r^2)_+^\alpha,
\]
the source derives uniform $L^p$ bounds for
$\alpha>\max\{0,n|1/p-1/2|-1/2\}$, and a maximal bound for finite
$p\ge2$ when $\alpha>s_+(p)$ \cite{consequences}.
The mechanism is worth explaining because it also prepares the square-function
argument.

Rescale the local-smoothing theorem to unit-frequency input and a long time
interval. For $\sigma>s_+(p)$,
\[
 \|U(t)g\|_{L^p_{x,t}(R\le|t|\le2R)}
 \lesssim R^{1/p+\sigma}\|g\|_p,
 \qquad R\ge1.
\]
A smooth frequency shell of relative width $\delta$ around radius $r$ has
an integral representation
\[
 \psi((1-|D|/r)/\delta)f
 =\frac{\delta r}{2\pi}\int e^{\mathrm itr}
                  \widehat\psi(\delta rt)U(t)g\dd t,
\]
where $g$ is an appropriate annular cutoff of the input. On one octave
$1\le r\le2$, take the supremum over $r$ and split the $t$ integral into
dyadic intervals. Holder in time gives a contribution bounded by
\[
 C_M\delta(1+\delta R)^{-M}R^{1+\sigma}\|g\|_p.
\]
Summing over $R$ costs $\delta^{-\sigma}$. Scale the argument over radius
octaves, and use Littlewood--Paley with $p\ge2$ to combine the inputs.

Near the sphere, the Bochner--Riesz multiplier is a sum of shell multipliers
weighted by $\delta^\alpha$. Choosing $s_+(p)<\sigma<\alpha$ makes
$\sum\delta^{\alpha-\sigma}$ converge. The part away from the sphere is a
smooth multiplier with integrable kernel. This proves the maximal estimate.
For each fixed radius, self-adjointness and duality extend the uniform estimate
to $1<p<2$. At $p=1,\infty$, the kernel is integrable for
$\alpha>(n-1)/2$, giving those endpoint bounds directly.

The distinction between the two types of uniformity matters:
$\sup_r\|B_r^\alpha f\|_p$ is not the same expression as
$\|\sup_r|B_r^\alpha f|\|_p$. The maximal argument proves the stronger second
expression only in its specified range. Norm convergence follows by density;
almost-everywhere convergence follows in the maximal range by approximation
and the maximal inequality.

\section{Stein's square function: keeping the radius integral inside}
\label{sec:square}
The square function is
\[
 \mathcal G_\alpha f(x)=
 \left(\int_0^\infty|r\partial_rB_r^\alpha f(x)|^2\frac{\dd r}{r}\right)^{1/2}.
\]
The source proves its $L^p$ bound for $2\le p<\infty$ and
\begin{equation}\label{eq:squarethreshold}
 \alpha>\frac12+s_+(p)
     =\max\left\{\frac12,n\left(\frac12-\frac1p\right)\right\}
\end{equation}
\cite{stein}. A supremum estimate alone does not establish this square
integral, since the radius measure has infinitely many scales.

On a single interval $1\le r\le2$, differentiating the multiplier near the
sphere costs $\delta^{\alpha-1}$. A shell of thickness $\delta$ yields
$\delta^{1/2}$ from the $L^2$ integral in radius. The wave estimate costs
$\delta^{-\sigma}$, where $\sigma>s_+(p)$. The product is
$\delta^{\alpha-1/2-\sigma}$, which is summable precisely in the strict range
\eqref{eq:squarethreshold}.

More explicitly, Plancherel in the radius variable transforms the shell
square norm into
\[
 C\delta\,
 \|\widehat\varphi(\delta t)U(t)g(x)\|_{L^2_t}.
\]
On a time interval of length comparable with $R$, Holder gives
$R^{1/2-1/p}$, and the rescaled wave estimate gives
$R^{1/p+\sigma}$. Summing dyadic $R$ thus yields a one-octave norm
$O(\delta^{1/2-\sigma})$. At $p=2$ the corresponding bound is
$O(\delta^{1/2})$ directly by spatial Plancherel.

To combine all radius octaves \emph{inside} the square norm, the source proves
a Hilbert-valued dilate lemma. On a test cube, very coarse scales are nearly
constant by a derivative estimate, and very fine scales are controlled by the
$L^2$ estimate plus the decaying kernel tail. Only
$O(1+\log(1/\delta))$ middle scales need the $L^p$ estimate. Holder over
those scales introduces the factor
$[1+\log(1/\delta)]^{1/2-1/p}$. A sharp-function inequality and the spatial
Littlewood--Paley estimate complete the vector-valued bound.

Consequently the full shell cost is bounded by
\[
 \delta^{\alpha-1}
 \left(\delta^{1/2}+\delta+
 [1+\log(1/\delta)]^{1/2-1/p}\delta^{1/2-\sigma}\right).
\]
Its dyadic sum converges because the inequality for $\alpha$ is strict.
The smooth part away from the sphere vanishes quadratically at zero and is
summed separately by ordinary multiplier bounds. At $p=2$ the exact calculation
is
\[
 \|\mathcal G_\alpha f\|_2^2
 =2\alpha^2\int_0^1u(1-u)^{2\alpha-2}\dd u\,\|f\|_2^2
 =\frac{\alpha}{2\alpha-1}\|f\|_2^2,
\]
valid for $\alpha>1/2$. This explicitly identifies the low-end square-function
threshold.

\section{Formal statements, proof routing, and verification limits}
\label{sec:formal}
Scope document 079 lists six targets: local smoothing, uniform Bochner--Riesz
bounds, spherical extension, the Kakeya maximal bound, full Hausdorff dimension,
and Stein's square function \cite{scope}. The reference file defines the actual
Fourier integrals, Bessel potentials, surface measure, tubes, and observation
interval, rather than using names for those objects without definitions.

The files under \texttt{ComparatorChallenges} are reference statements. Their
proof placeholders are not the solution proof. The JSON configuration routes
the targets to the separate solution module
\texttt{OAI.Analysis.Fourier.LocalSmoothing.Comparator} and specifies the
permitted foundational axioms \texttt{propext}, \texttt{Quot.sound}, and
\texttt{Classical.choice} \cite{challenge,config,solution}.
The presence of this configuration is evidence about the intended verification
procedure; it is not a record of a successful independent run by this preparer.

The main distinction is between three levels of assertion. First, the
manuscript supplies the natural-language analytic proof reconstructed here.
Second, the scope document and comparator identify the reported formal end
statements. Third, this report supplies explanation and catalogue matching.
It does not silently turn the third level into a new certification of the
first two.

In particular, a formal target for every positive Sobolev loss does not, by
itself, certify the precise logarithmic annular exponent, every auxiliary
barrier, or every refinement in the manuscript. These are separate statements
whose exact formal routing would have to be checked. The development was not
compiled here, no kernel or comparator run is claimed, and no independent
expert endorsement is inferred from the existence of public source files.

\section{What has been expanded and what a complete audit would require}
\label{sec:audit}
Compared with the shorter report, this edition supplies the actual chain from
the first-order symbol to the positive current, the role of the holomorphic
lift and first integral, the determinant/cofactor calculus, the moving-fiber
Dirichlet argument, the corrected-energy cancellation, the paired barriers,
and the two radial comparisons. It also supplies the main deduction mechanisms
for the four families of consequences, rather than only listing their names.

The most consequential independent analytic checks are still the transformed
current's conservation, the full fiber geometry and boundary compression,
the Dirichlet a priori estimates, the moving-domain energy identity, and the
uniform scale dependence of the barriers. These are the points at which a
formal-looking calculation can fail if a sign, a boundary term, or a
frequency-dependent constant is overlooked. The source proofs, not the shorter
summaries in this report, remain the controlling reference for those details.

A full formal audit would rebuild the pinned solution dependency closure, run
the configured comparison, inspect the accepted axioms, and compare the final
definitions and quantifiers with Math-273. None of these full audits is claimed
as a completed task here. This report also does not attempt to reproduce every
appendix on angular regularity, generalized pencils, planar variants, or
parabolic consequences; those extensions are not required for its primary
catalogue claim.

\section{Completion Estimate}
\noindent\textbf{COMPLETION: 0\%}
This report claims no new original solution or independently verified proof. The estimate concerns this report's original research contribution, not the extent of the results attributed to OpenAI.
\section{Final status and preparation}
Unresolved as an original research attempt: literature exposition only. OpenAI's reported full result remains attributed to the cited source; no source-catalogue or community-review verdict is changed by this report.
The document is dated 26 October 2026. The underlying GPT 6 Astra Pro release/build and reasoning setting are unknown. The supplied source-inspection statements describe the preparation of the report; repository integration checks identity, formatting, compilation and rendering, rather than certifying the underlying proofs. Original generation prompts and browsing logs are unavailable; use of the repository's initial and final research prompts during original preparation is not documented. No independent human audit, Lean execution or new performance benchmark is asserted.


\section{Attribution and provenance}
\paragraph{Problem handle.} \texttt{prb\_94f7553139864cf78d6311e3db49c212}
(Math-273).
\paragraph{Submission claim.} Literature report.
\paragraph{Underlying reported result.} Full solution of Math-273, attributed to
OpenAI, with the documented formalization scope described above.
\paragraph{Exact scope.} A detailed source-pinned exposition of the proof's main
analytic constructions and deductions. It explains the full-dimensional Kakeya
consequence and preserves the strict-boundary qualification for the inspected
restriction statement. It claims neither a new discovery nor an independently
verified complete proof.
\paragraph{Who did the work.} LLM only (people only prompted or submitted).
\paragraph{Model used in this preparation.} OpenAI GPT 6 Astra Pro; exact version:
unknown; thinking setting: unknown; actual time or compute budget: unknown.
Role: source inspection, mathematical exposition, scope matching, and TeX
preparation. No second model or independent human proof reviewer was used.
\paragraph{People and roles.} The submitter prompted preparation and expansion. No human mathematical authorship or independent proof-verification role is asserted.
\paragraph{External research provenance.} The manuscript names OpenAI as author.
The internal models, settings, budgets, and individual roles behind the original
research are not established here and are recorded as unknown. They are source
provenance, not additional models invoked to write this report.
Linked third-party artifacts retain their original terms. Source manuscripts, images and Lean code are not bundled.

\begin{thebibliography}{99}
\bibitem{catalogue} \emph{we-solve}, Math-273, \emph{Local Smoothing for the
Euclidean Wave Equation}. Used: problem identity, target, and earlier
three-dimensional entry. \url{https://we-solve.github.io/problems/local-smoothing-for-the-euclidean-wave-equation/}.
\bibitem{intro} OpenAI, \emph{Local smoothing for the Euclidean wave equation},
October 25, 2026, pinned commit stated in Section~\ref{sec:target}.
\href{https://github.com/openai/math/blob/7ef5c0cc5ca4b0e072b7a050fca09ab24bc6fa3a/preprints/Local-smoothing-for-the-Euclidean-wave-equation-October-25-2026/build/sections/introduction.tex}{Introduction and main theorem statements}.
Used: local-smoothing and critical annular statements and the proof roadmap.
\bibitem{current} OpenAI, same manuscript,
\href{https://github.com/openai/math/blob/7ef5c0cc5ca4b0e072b7a050fca09ab24bc6fa3a/preprints/Local-smoothing-for-the-Euclidean-wave-equation-October-25-2026/build/sections/current.tex}{Conserved-current construction}.
Used: first-order system, analytic continuation, paired-state current,
holomorphic lift, first integral, and scalar window detection.
\bibitem{calculus} OpenAI, same manuscript,
\href{https://github.com/openai/math/blob/7ef5c0cc5ca4b0e072b7a050fca09ab24bc6fa3a/preprints/Local-smoothing-for-the-Euclidean-wave-equation-October-25-2026/build/sections/calculus.tex}{Quaternionic Hessian and cofactor calculus}.
Used: product, divergence, Schur-complement, and normal-compression identities.
\bibitem{energy} OpenAI, same manuscript,
\href{https://github.com/openai/math/blob/7ef5c0cc5ca4b0e072b7a050fca09ab24bc6fa3a/preprints/Local-smoothing-for-the-Euclidean-wave-equation-October-25-2026/build/sections/energy.tex}{Moving fibers and corrected energy}.
Used: fiber geometry, rank-one background, regularization, determinant equation,
energy balance and energy comparison.
\bibitem{dirichlet} OpenAI, same manuscript,
\href{https://github.com/openai/math/blob/7ef5c0cc5ca4b0e072b7a050fca09ab24bc6fa3a/preprints/Local-smoothing-for-the-Euclidean-wave-equation-October-25-2026/build/sections/dirichlet.tex}{Fiber Dirichlet theorem}.
Used: existence, comparison, boundary and interior a priori estimates,
regularity, and smooth dependence on parameters.
\bibitem{barrier} OpenAI, same manuscript,
\href{https://github.com/openai/math/blob/7ef5c0cc5ca4b0e072b7a050fca09ab24bc6fa3a/preprints/Local-smoothing-for-the-Euclidean-wave-equation-October-25-2026/build/sections/barrier.tex}{Boundary and interior barriers}.
Used: flux decay, interior detection, order of parameter choices, and the scalar
differential inequality.
\bibitem{propagation} OpenAI, same manuscript,
\href{https://github.com/openai/math/blob/7ef5c0cc5ca4b0e072b7a050fca09ab24bc6fa3a/preprints/Local-smoothing-for-the-Euclidean-wave-equation-October-25-2026/build/sections/propagation.tex}{Scalar propagation between circles}.
Used: radial comparison, physical-time transfer, and the two logarithmic losses.
\bibitem{transfer} OpenAI, same manuscript,
\href{https://github.com/openai/math/blob/7ef5c0cc5ca4b0e072b7a050fca09ab24bc6fa3a/preprints/Local-smoothing-for-the-Euclidean-wave-equation-October-25-2026/build/sections/transfer.tex}{Passage to Euclidean space and Sobolev estimates}.
Used: cap summation, period removal, interpolation and ordinary/logarithmic
Sobolev summation.
\bibitem{sharpness} OpenAI, same manuscript,
\href{https://github.com/openai/math/blob/7ef5c0cc5ca4b0e072b7a050fca09ab24bc6fa3a/preprints/Local-smoothing-for-the-Euclidean-wave-equation-October-25-2026/build/sections/sharpness.tex}{Endpoint obstructions}.
Used: the radial focusing construction and its stationary-phase lower bound.
The additional weak-norm construction is not reproduced here.
\bibitem{kakeya-proof} OpenAI, same manuscript,
\href{https://github.com/openai/math/blob/7ef5c0cc5ca4b0e072b7a050fca09ab24bc6fa3a/preprints/Local-smoothing-for-the-Euclidean-wave-equation-October-25-2026/build/sections/kakeya.tex}{Kakeya and local extension arguments}.
Used: focusing-shell extension, weighted tube overlap, maximal estimates,
and the multiscale Hausdorff argument.
\bibitem{global-extension} OpenAI, same manuscript,
\href{https://github.com/openai/math/blob/7ef5c0cc5ca4b0e072b7a050fca09ab24bc6fa3a/preprints/Local-smoothing-for-the-Euclidean-wave-equation-October-25-2026/build/sections/global-extension.tex}{Global extension}.
Used: the source's application of epsilon removal, interpolation, and slicing.
The earlier epsilon-removal theorem is credited through this inspected source;
its original proof was not independently audited here.
\bibitem{consequences} OpenAI, same manuscript,
\href{https://github.com/openai/math/blob/7ef5c0cc5ca4b0e072b7a050fca09ab24bc6fa3a/preprints/Local-smoothing-for-the-Euclidean-wave-equation-October-25-2026/build/sections/consequences.tex}{Bochner--Riesz means}.
Used: long-time wave rescaling, shell representation, maximal summation,
duality, and convergence consequences.
\bibitem{stein} OpenAI, same manuscript,
\href{https://github.com/openai/math/blob/7ef5c0cc5ca4b0e072b7a050fca09ab24bc6fa3a/preprints/Local-smoothing-for-the-Euclidean-wave-equation-October-25-2026/build/sections/stein_square.tex}{Stein's square function}.
Used: one-octave shell estimate, Hilbert-valued dilate lemma, all-scale summation,
and the exact $L^2$ identity.
\bibitem{scope} OpenAI,
\href{https://github.com/openai/math/blob/7ef5c0cc5ca4b0e072b7a050fca09ab24bc6fa3a/lean/docs/079.md}{Formalization scope 079}.
Used: the exact documented end statements and their parameter ranges.
\bibitem{challenge} OpenAI,
\href{https://github.com/openai/math/blob/7ef5c0cc5ca4b0e072b7a050fca09ab24bc6fa3a/lean/ComparatorChallenges/LocalSmoothing.lean}{Local-smoothing comparator reference}.
Used: target definitions and signatures; this is not the solution proof.
\bibitem{config} OpenAI,
\href{https://github.com/openai/math/blob/7ef5c0cc5ca4b0e072b7a050fca09ab24bc6fa3a/lean/ComparatorChallenges/LocalSmoothing.json}{Local-smoothing comparator configuration}.
Used: module routing and permitted axioms.
\bibitem{solution} OpenAI,
\href{https://github.com/openai/math/blob/7ef5c0cc5ca4b0e072b7a050fca09ab24bc6fa3a/lean/OAI/Analysis/Fourier/LocalSmoothing/Comparator.lean}{Local-smoothing solution module}.
Used: header, definitions and routing; no complete dependency audit or execution
is claimed.
\bibitem{kakeya-catalogue} \emph{we-solve}, Math-95,
\emph{Kakeya conjecture in dimensions four and higher}.
Used: Hausdorff-dimension scope comparison.
\url{https://we-solve.github.io/problems/kakeya-conjecture-in-dimensions-four-and-higher/}.
\bibitem{restriction-catalogue} \emph{we-solve}, Math-96,
\emph{Sharp exponents in the Fourier restriction conjecture}.
Used: the distinction between its boundary formulation and the inspected
strict-interior theorem.
\url{https://we-solve.github.io/problems/sharp-exponents-in-the-fourier-restriction-conjecture/}.
\end{thebibliography}
\end{document}
